% presentation/slides/11_conclusion.tex
\begin{frame}{Conclusion et Perspectives}
    \textbf{Trois enseignements fondamentaux :}
    \begin{enumerate}
        \item \textbf{Le coût statistique du masque :} L'absence géométrique d'information ($Z_i$ cachés) détruit la linéarité classique et introduit de la non-convexité via le profil \textit{log-sum-exp}.
        \item \textbf{La puissance découplée de l'EM :} L'algorithme EM s'avère optimal en pratique en reformulant l'étape M complexe en itérations de sous-problèmes convexes et linéaires stables.
        \item \textbf{Complémentarité Statistique-Régularisation :}} La quantité de données ($n \to \infty$) réduit l'erreur de variance, mais seule la pénalisation mathématique structurelle gère le caractère mal posé de la physique (Radon 2D).
    \end{enumerate}
    
    \vspace{0.2cm}
    \textbf{Ouverture :} Généralisation en 3D (espace des rotations non commutatif $SO(3)$) et traitement des hétérogénéités de distribution de $f_Z$.
\end{frame}